\chapter{数据分析}
\section{数据清洗与转换 (dplyr \& tidyr)}

在 R 语言的实际应用中，原始数据往往是“脏”的或格式不规范的。\texttt{tidyverse} 生态系统中的 \texttt{dplyr} 和 \texttt{tidyr} 提供了一套高度一致、语义化的工具。

\subsection{dplyr 的核心动作 (Verbs)}
\texttt{dplyr} 包通过几个核心动词函数，覆盖了大部分数据转换需求：

\begin{enumerate}[label=\arabic*., leftmargin=2em]
    \item \textbf{filter()}：按行筛选数据。
    \item \textbf{select()}：按列选取变量。
    \item \textbf{mutate()}：衍生新列。
    \item \textbf{arrange()}：数据排序。
    % 修正点：将 group_by 写在 \texttt{} 中并转义下划线
    \item \textbf{\texttt{group\_by()}} 与 \textbf{\texttt{summarise()}}：分组聚合。
\end{enumerate}

\subsection{tidyr 的数据重塑}
% 修正点：对 pivot_longer 和 pivot_wider 进行转义
\begin{itemize}
    \item \texttt{pivot\_longer()}：将宽格式数据转换为长格式。
    \item \texttt{pivot\_wider()}：将长格式数据转换为宽格式。
\end{itemize}



\subsection{常用清洗函数对比总结}
\begin{table}[h!]
\centering
\renewcommand{\arraystretch}{1.5}
\begin{tabular}{|l|p{6cm}|l|}
\hline
\textbf{函数} & \textbf{功能描述} & \textbf{SQL/Excel 对应} \\ \hline
\texttt{filter()} & 筛选满足条件的观察行 & Filter / WHERE \\ \hline
\texttt{select()} & 选取或剔除指定的变量列 & Select / Column hide \\ \hline
\texttt{mutate()} & 新增列并保留原数据 & Calculated Field \\ \hline
% 修正点：表格中的下划线处理
\texttt{group\_by()} & 将数据按类别分组 & GROUP BY \\ \hline
\texttt{summarise()} & 将多行数据压缩为一行汇总结果 & Pivot Table / Sum \\ \hline
\texttt{inner\_join()} & 根据键值合并两个数据框 & VLOOKUP / JOIN \\ \hline
\end{tabular}
\end{table}





\section{数据可视化 (ggplot2)}

\texttt{ggplot2} 是 R 语言中最强大的可视化包。其核心逻辑是通过加号 (\texttt{+}) 将不同的图层叠加在一起。

\subsection{ggplot2 的基础语法结构}
一个基本的绘图模板如下：
\begin{lstlisting}[language=R]
ggplot(data = <DATA>) + 
  <GEOM_FUNCTION>(mapping = aes(<MAPPINGS>)) +
  <THEME_FUNCTION>()
\end{lstlisting}

\begin{enumerate}[label=\arabic*., leftmargin=2em]
    \item \textbf{数据 (Data)}：必须是一个数据框 (\texttt{data.frame})。
    \item \textbf{映射 (Mapping)}：使用 \texttt{aes()} 函数定义，指定变量如何映射到视觉属性（如 \texttt{x} 轴、\texttt{y} 轴、颜色 \texttt{color}、形状 \texttt{shape}）。
    \item \textbf{几何对象 (Geoms)}：决定图表类型（点、线、柱等）。
\end{enumerate}



\subsection{常见图表类型 (Geoms)}
根据研究目的选择不同的几何对象函数：

\begin{itemize}
    \item \textbf{散点图}：\texttt{geom\_point()} —— 用于观察两个连续变量的关系。
    \item \textbf{箱线图}：\texttt{geom\_boxplot()} —— 用于观察变量的分布及异常值。
    \item \textbf{条形图}：\texttt{geom\_bar()} 或 \texttt{geom\_col()} —— 用于展示分类数据的频数或数值。
    \item \textbf{折线图}：\texttt{geom\_line()} —— 常用于时间序列数据。
    \item \textbf{直方图}：\texttt{geom\_histogram()} —— 观察单变量的分布密度。
\end{itemize}

\subsection{图形美化：颜色、标签与主题}
通过叠加额外的图层，可以让图表达到出版标准。
\begin{lstlisting}[language=R]
df %>%
  ggplot(aes(x = weight, y = height, color = group)) +
  geom_point(size = 3, alpha = 0.7) +
  labs(
    title = "Weight vs Height",
    x = "Weight (kg)",
    y = "Height (cm)",
    color = "Group Name"
  ) +
  theme_bw() # 使用经典白底黑框主题
\end{lstlisting}

\subsection{分面 (Faceting)}
分面是 \texttt{ggplot2} 最强大的功能之一，可以根据某个分类变量快速生成多张子图。
\begin{itemize}
    \item \texttt{facet\_wrap(\~{}variable)}：按一个变量自动排列子图。
    \item \texttt{facet\_grid(rows \~{} cols)}：按行和列两个变量生成矩阵子图。
\end{itemize}



\subsection{常用外观调整函数对照表}
\begin{table}[h!]
\centering
\renewcommand{\arraystretch}{1.5}
\begin{tabular}{|l|l|p{6cm}|}
\hline
\textbf{类别} & \textbf{函数} & \textbf{功能说明} \\ \hline
\texttt{坐标轴范围} & \texttt{xlim(), ylim()} & 设置 X 或 Y 轴的起止点 \\ \hline
\texttt{颜色比例} & \texttt{scale\_color\_manual()} & 手动指定离散变量的颜色（如实验组红色，对照组蓝色） \\ \hline
\texttt{坐标翻转} & \texttt{coord\_flip()} & 将 X 轴与 Y 轴对调（常用于条形图） \\ \hline
\texttt{预设主题} & \texttt{theme\_minimal()} & 极简主题，去除冗余网格线 \\ \hline
\texttt{保存图表} & \texttt{ggsave()} & 将最后生成的图表保存为 pdf, png 或 tiff 格式 \\ \hline
\end{tabular}
\end{table}

\paragraph{注意事项：}
1. 注意 \texttt{color}（轮廓色/点色）与 \texttt{fill}（填充色）的区别。柱状图的颜色通常使用 \texttt{fill}。
2. 多个图层之间必须使用 \texttt{+} 连接，且 \texttt{+} 必须写在行末，不能写在行首。
3. 如果在 \texttt{aes()} 内部设置颜色，它是根据数据自动映射；如果在 \texttt{aes()} 外部设置颜色（如 \texttt{geom\_point(color = "blue")}），则是统一指定颜色。



\section{统计检验基础}

R 语言提供了强大的内置函数，用于执行常见的参数化与非参数化检验。在执行检验前，通常需要确保数据满足正态性（\texttt{shapiro.test()}）和方差齐性（\texttt{bartlett.test()}）。

\subsection{均值比较 (T-Test \& ANOVA)}

\begin{enumerate}[label=\arabic*., leftmargin=2em]
    \item \textbf{T 检验 (t-test)}
    用于比较两组数据均值是否存在显著差异。
    \begin{itemize}
        \item \textbf{独立样本 T 检验}：比较两组独立受试者。
        \item \textbf{配对样本 T 检验}：比较同一受试者治疗前后的变化。
    \end{itemize}
\begin{lstlisting}[language=R]
# 独立样本：y 是连续变量，group 是二分类因子
t.test(y ~ group, data = df)

# 配对样本
t.test(df$before, df$after, paired = TRUE)
\end{lstlisting}

    \item \textbf{方差分析 (ANOVA)}
    用于比较三组及以上数据的均值差异。
\begin{lstlisting}[language=R]
# 1. 拟合线性模型
fit <- aov(weight ~ group, data = df)

# 2. 查看 ANOVA 表
summary(fit)

# 3. 事后两两比较 (Tukey HSD)
TukeyHSD(fit)
\end{lstlisting}
\end{enumerate}



\subsection{相关性与回归分析}

\begin{enumerate}[label=\arabic*., leftmargin=2em]
    \item \textbf{相关性分析}
    计算变量间的相关系数（Pearson 或 Spearman）。
\begin{lstlisting}[language=R]
cor(df$age, df$blood_pressure, method = "pearson")
cor.test(df$age, df$blood_pressure) # 包含显著性检验
\end{lstlisting}

    \item \textbf{线性回归 (Linear Regression)}
    用于探究自变量 ($x$) 对因变量 ($y$) 的影响。
\begin{lstlisting}[language=R]
# 拟合模型：y = ax + b
model <- lm(y ~ x1 + x2, data = df)

# 查看结果（包括系数、R方和 P 值）
summary(model)

# 提取残差
residuals(model)
\end{lstlisting}
\end{enumerate}



\subsection{常用统计函数速查表}
\begin{table}[h!]
\centering
\renewcommand{\arraystretch}{1.5}
\begin{tabular}{|l|l|p{6cm}|}
\hline
\textbf{统计需求} & \textbf{R 函数} & \textbf{适用场景} \\ \hline
\texttt{正态性检验} & \texttt{shapiro.test()} & 判断数据是否符合正态分布 ($P > 0.05$) \\ \hline
\texttt{二组均值} & \texttt{t.test()} & 两组连续变量比较 \\ \hline
\texttt{二组中位数} & \texttt{wilcox.test()} & 非正态分布数据的两组比较 \\ \hline
\texttt{多组均值} & \texttt{aov()} & 三组及以上均值比较 (ANOVA) \\ \hline
\texttt{分类变量检验} & \texttt{chisq.test()} & 观察频数与期望频数的差异 (卡方检验) \\ \hline
\texttt{线性模型} & \texttt{lm()} & 回归分析与预测 \\ \hline
\end{tabular}
\end{table}

\paragraph{注意事项：关于 P 值}
1. 在 R 的输出结果中，星号通常代表显著性水平：\texttt{***} ($P < 0.001$), \texttt{**} ($P < 0.01$), \texttt{*} ($P < 0.05$), \texttt{.} ($P < 0.1$)。
2. 对于回归模型，在使用 \texttt{summary(model)} 后，重点关注 \texttt{Coefficients} 表中的 \texttt{Pr(>|t|)} 列以及底部的 \texttt{Adjusted R-squared}（调整后 R 方）。